\documentclass[12pt, a4paper]{article}

% --- Language and Font Encoding ---
\usepackage[utf8]{inputenc}
\usepackage[T1]{fontenc}
\usepackage[czech]{babel}

% --- Elegant Font (Palatino) ---
\usepackage{tgpagella}

% --- Page Geometry (The A4 Secret) ---
% We use wide margins so the poem doesn't look lost on the large A4 sheet.
% This creates a beautiful, professional "frame" of white space.
\usepackage[left=4.5cm, right=4.5cm, top=4cm, bottom=4.5cm]{geometry}

% --- Spacing and Poetry Formatting ---
\usepackage{setspace}
\usepackage{verse}

% --- Slightly Fancy Section Titles ---
\usepackage{titlesec}
\titleformat{\section}[block]{\Large\bfseries\itshape\centering}{}{0pt}{}

% --- Subtle Ornament to separate sections ---
\newcommand{\ornament}{
    \vspace{1.5em}
    \begin{center}
        $\ast$\quad$\ast$\quad$\ast$
    \end{center}
    \vspace{1.5em}
}

% --- Document Title Setup ---
\title{\Huge\bfseries\itshape Proč neužít dne?}
\author{Dominik Dembinný} % You can type your name inside these brackets!
\date{}   % This hides the date. Remove this line if you want the current date to show.

\begin{document}


\maketitle
\thispagestyle{empty} 
\vspace{2em} % Adds some breathing room below the title before the poem starts

% 1.5 line spacing gives the text an airy, poetic feel
\onehalfspacing

% ================= SECTION 1 =================
\section*{Část I.}

% The verse environment centers the block of text on the page 
% while keeping the lines themselves left-aligned.
\begin{verse}[0.85\textwidth]

Ptáček se vznáší a zkoumá ten svět, \\
křídly se dotýká nebeských drah. \\
Hledá tu pravdu a touží hned vzlét, \\
zmizel už navěky pozemský strach.

\vspace{1em}

Myslí si bláhově: láska je mír, \\
spojení duší jak hedvábný sen. \\
V hladině rybníka utichl vír, \\
sladce se objímat celičký den. 

\vspace{1em}

Láska není rybník snů, \\
nechce ticho, nechce klid! \\
Je to válka tvůrců dnů, \\
musíš staré dlátem bít!

\vspace{1em}

Zmatený ptáček teď zkouší dál žít, \\
k bližnímu upíná zraněný zrak. \\
Chce jenom bezpečí, teplo a klid, \\
schovat se do hnízda, odletět v mrak.

\vspace{1em}
\newpage

Slabosti omluvit, smutky si třít, \\
milovat toho, kdo blízko teď stál. \\
V náruči hřejivé slaboučký být, \\
aby se nikdo už bolesti bál.

\vspace{1em}

Pryč s tou láskou k těm, co znáš, \\
spal svůj starý, sláblý svět! \\
Vyšší bytost v ohni máš, \\
z ohně musí vzlétnout hned!

\vspace{1em}

Zděšeně pohlíží na ostrý svit, \\
nad propast hlubokou upírá zrak. \\
Vidí tam zvířata v objetí dlít, \\
hada a orla, co letí jak mrak.

\vspace{1em}

Myslí, že hadek se v peříčku hřál, \\
v hnízdečku teplém a v tichosti spí. \\
Orel, ten dobrák, se z výšiny smál, \\
o lásce něžné a mírné teď sní.

\vspace{1em}

Žádné hnízdo, žádný klid, \\
orel nese hada výš. \\
V tenzi musí oba bdít, \\
jinak spadnou v černou skrýš!

\vspace{1em}

Ptáček už chápe, co znamená moc, \\
rozbíjet sebe a vznášet se dál. \\
Protančit v bolesti celičkou noc, \\
aby se nadčlověk v záři teď smál.

\vspace{1em}
\newpage

Rozechvěně vyletuje, \\
veseleji pokračuje. \\
Usměje se do hluboka, \\
zapomene na proroka. \\
Opaduje unaveně, \\
rozplyne se do plamene.

\end{verse}


\ornament

% ================= SECTION 2 =================
\section*{Část II.}

\begin{verse}[0.85\textwidth]

Proč se musím lámat v půli, \\
když je snazší klidně žít? \\
Pít a zpívat podle vůle, \\
v prázdnu ale chybí cit.

\vspace{1em}

Mám se tedy svázat řádem,\\
plnit jenom přísný mrav? \\
Zákon ale duši svírá, \\
nenajdu v něm pravý stav.

\vspace{1em}

Ta úzkost není slabost, nýbrž brána, \\
jíž musím projít, abych našel směr. \\
Já neschovám se jako líná vrána, \\
chci čelit prázdnu, překročit svůj stín.

\vspace{1em}

Však nadčlověk je příliš pyšná maska, \\
vždyť pouhou vůlí nezkrotím ten svět. \\
Když selže rozum, začne pravá láska, \\
a nemožné se stane pravdou hned.

\vspace{1em}

Rozum mě svazuje těžkými lany, \\
víra však žádá ten absurdní skok. \\
Pochybnost pálí jak do živé rány, \\
zastavit nechci ten bláznivý tok. \\
Tam, kde už končí mé bezpečné kroky, \\
začíná věčnost a smazává rok.

\vspace{1em}


Pravda se neskrývá v jásavém davu, \\
musí se zrodit a najít tvůj cíl. \\
Ztrať se teď v sobě a skloň pyšnou hlavu, \\
v tichosti nitra je největší díl. \\
Vášeň a niternost, to je ta spása, \\
Bohu když odevzdáš celičký svět.

\vspace{1em}

Rytíř tvé víry se navenek směje, \\
vypadá prostě a tiše tu jde. \\
Uvnitř však nekonečno se v něm chvěje, \\
vzdal se už všeho a získal vše zpět. \\
Není to nadčlověk s kladivem v ruce, \\
obyčejností svou dobývá svět.

\vspace{1em}

Přestal jsem bořit a kladivo padá, \\
nacházím sebe, ač ztrácím svou moc. \\
Duše už neválčí, klidně se skládá, \\
ve věčné síle se rozpustí noc. \\
To, co mě stvořilo, drží mé bytí, \\
konečně zmizel ten úzkostný strach.

\vspace{1em}
\newpage

Pomaličku odhazuji tíhu, \\
do věčnosti dál. \\
Odevzdávám roztříštěnou duši, \\
prázdný je můj štít. \\
Nehledám už nadpozemskou sílu, \\ 
nechci o ni stát. \\
Nacházím své opravdové jádro, \\
můžu se teď smát.

\vspace{1em}


Pomaličku rozplývám se v Bohu, \\
ztrácí se mi klam. \\
Odpočívám, nehledám už cestu, \\
nejsem na to sám. \\
Nekonečno potkává se s časem, \\
padám do těch vln. \\
Věřím v to, co přesahuje rozum, \\
naprosto jsem pln.

\end{verse}

\ornament

% ================= SECTION 3 =================
\section*{Část III.}

\begin{verse}[0.85\textwidth]

Proč neužít života, neužít dne? \\
Když slunce se do duše laskavě pne, \\
a úzkost, co drtila, rázem se hne. \\
Už netlačí břemeno, netíží kříž, \\
už nepadám do prázdna, stoupám teď výš, \\
to ticho teď hojí a dotýká mne.

\vspace{1em}
\newpage

Proč neužít života, neužít dne? \\
Když víno teď v pohárech jiskří a žhne, \\
a radost se v každičké buňce teď dme. \\
Jen opustit logiku, vnímat ten proud, \\
a nechat se v extázi do dálky plout, \\
ta chvíle se k našemu osudu lne.

\vspace{1em}


Proč neužít života, neužít dne? \\
Když hudba teď do smyslů divoce tne, \\
a kůže se druhého zlehýnka tkne. \\
V tom víru se ztrácejí hranice těl, \\
kdo před chvílí plakal, ten rázem by pěl, \\
a potůček potu už pomalu schne.

\vspace{1em}

Proč neužít života, neužít dne? \\
Když barva se rozpadá, voní a žhne, \\
a obrys se v mlhavém prostoru hne. \\
Už nevnímáš včera a nevnímáš svět, \\
jsi nektar a omamný, zářivý květ, \\
a krása mě pohltí, celého mne.

\vspace{1em}

Proč neužít života, neužít dne? \\
Když paprsek světla se temnoty tkne, \\
a prostor se do křídel široce pne. \\
My plujeme na vlnách tekutých chvil, \\
už neznáme cesty a nemáme cíl, \\
a věčnost se k zástupům hvězdokup lne.

\vspace{1em}
\newpage

Proč neužít života, neužít dne? \\
Když iluze propadá, do uší zne, \\
a časová smyčka se trhavě hne. \\
Vše pulzuje, praská a stéká jak med, \\
to není už hmota a není to led, \\
to duše se do šířky bláznivě dme.

\vspace{0.5em}

Proč neužít života, neužít dne? \\
Když ozvěna krystalů do očí žhne, \\
a dech už se zastavil, stěží se hne. \\
Jsi vibrace atomů, jiskřivý prach, \\
už dávno se rozpustil poslední strach, \\
a zlato se do masa pomalu tne.

\vspace{0.5em}

Proč neužít života, neužít dne? \\
Ta vůně se sytí a do prstů lne, \\
ty stíny se barví a ticho se dme. \\
Jen hluboké tóny a zářivý šum, \\
co boří ten znavený, falešný dům, \\
a konečná pravda se přede mnou pne.

\vspace{0.5em}

Proč neužít života, neužít dne? \\
Když vrtise krtise do mraků žhne, \\
a trnkatá zvukohra oblakem hne. \\
Ty třepoty kmitoty barvitý let, \\
ten modravý mihotav blaživý svět, \\
a hvizdota laskota k oku se lne.

\vspace{0.5em}

Proč neužít života, neužít dne? \\
Když tatata lalala v plutu teď fle, \\
a báboly čáraluk k aleku shne. \\
an lákala zárula ručí ku pál, \\
rak trmaca grmaca vígr dy vza, \\
la hula op capadop, tak užívej dne.

\end{verse}

\end{document}